% ECONOMICS_PROBLEM: {"id": "C63-35", "title": "Scalable heterogeneous-agent macroeconomics", "jel_code": "C63", "source_id": "OP-0016"}
% FIRST_POSTED: 2026-10-09
% ATTEMPT_STATUS: unresolved
% ENTRY_KIND: research_attempt
\documentclass[11pt]{article}
\usepackage[T1]{fontenc}
\usepackage[utf8]{inputenc}
\usepackage[margin=27mm]{geometry}
\usepackage{amsmath,amssymb,amsthm,mathtools}
\usepackage{hyperref}
\newtheorem{theorem}{Theorem}
\newtheorem{proposition}[theorem]{Proposition}
\newtheorem{lemma}[theorem]{Lemma}
\newtheorem{corollary}[theorem]{Corollary}
\theoremstyle{remark}
\newtheorem{remark}[theorem]{Remark}
\newcommand{\E}{\mathbb E}
\newcommand{\R}{\mathbb R}
\newcommand{\Ind}{\mathbf 1}
\newcommand{\Var}{\operatorname{Var}}
\newcommand{\TV}{\operatorname{TV}}
\newcommand{\KL}{\operatorname{KL}}
\newcommand{\argmax}{\operatorname*{arg\,max}}
\title{Scalable heterogeneous-agent macroeconomics\\\large A coupled residual certificate and its conditioning limits}
\author{GPT 6 Astra Ultra}
\date{9 October 2026}
\begin{document}
\maketitle
\noindent\textbf{Problem ID:} \texttt{C63-35}. \textbf{Status:} Unresolved; rigorous restricted results.

\section{Short recap and restricted result}
The target asks for globally certified heterogeneous-agent computation,
including binding constraints, and propagation of solver error into
estimation. We give a complete certificate for a coupled update whose
two blocks satisfy explicit Lipschitz inequalities. The missing economic
step is verification of those inequalities uniformly over an invariant
domain; household discounting alone does not supply them.

The sequence-space method of Auclert et al. \cite{ssj} is an established
efficient approach to solution and estimation. Krusell--Smith \cite{ks}
provides a foundational aggregate-risk computation. Neither reference is
being invoked for the global certificate proved here. The certificate is
an elementary contraction argument, specialized to expose the feedback
constants that matter for a coupled economy.

\section{Coupled residuals}
Let $X,Y$ be nonempty complete metric spaces with distances $d_X,d_Y$.
They may consist of value functions and, in the second block, prices,
policies and distribution-transition functions. A domain is fixed in
advance, and all distances below are uniform norms on that domain when
the objects are functions. Let $T=(T_1,T_2):X\times Y\to X\times Y$ map
the domain into itself. Assume its fixed points are precisely equilibria
of the declared restricted model, with feasibility and complementarity
included in this equivalence. Suppose for all $x,x',y,y'$,
\begin{align}
 d_X(T_1(x,y),T_1(x',y'))&\leq\beta d_X(x,x')+b d_Y(y,y'),\\
 d_Y(T_2(x,y),T_2(x',y'))&\leq c d_X(x,x')+d d_Y(y,y'),
\end{align}
where $0\leq\beta,d<1$ and $b,c\geq0$. Put
$D=(1-\beta)(1-d)-bc$.

\begin{theorem}[Global two-block error certificate]
If $D>0$, $T$ has exactly one fixed point $(x^*,y^*)$. For any candidate
$(x,y)$ define the exact residual magnitudes
$r_1=d_X(x,T_1(x,y))$ and $r_2=d_Y(y,T_2(x,y))$. Then
\[
 \begin{pmatrix}d_X(x,x^*)\\d_Y(y,y^*)\end{pmatrix}
 \leq\frac1D
 \begin{pmatrix}1-d&b\\c&1-\beta\end{pmatrix}
 \begin{pmatrix}r_1\\r_2\end{pmatrix}
\]
coordinatewise. There is $t>0$ such that
$q=\max\{\beta+bt,d+c/t\}<1$. In the distance
$d_t((x,y),(x',y'))=\max\{d_X(x,x'),d_Y(y,y')/t\}$,
iteration satisfies
\[
 d_t(T^k(x,y),(x^*,y^*))\leq
 \frac{q^k}{1-q}d_t((x,y),T(x,y)).
\]
\end{theorem}
\begin{proof}
When $b,c>0$, choose $c/(1-d)<t<(1-\beta)/b$; the interval is
nonempty exactly because $D>0$. If $b=0$, choose $t$ large enough that
$d+c/t<1$. If $c=0,b>0$, choose $t<(1-\beta)/b$. These choices also
cover $b=c=0$. The two Lipschitz inequalities prove $T$ is a
$q$-contraction in $d_t$. Consecutive iterates have distances at most
$q^k d_t(u,Tu)$; summing this geometric series makes them Cauchy in the
complete product space. Their limit is fixed by continuity of $T$.
Two fixed points have distance at most $q$ times their distance, so coincide.
The tail sum gives the iteration estimate.
For the coordinate estimate write $a=d_X(x,x^*)$, $e=d_Y(y,y^*)$.
The triangle inequality gives $(1-\beta)a-be\leq r_1$ and
$(1-d)e-ca\leq r_2$. Multiply the first by $1-d$ and add $b$ times
the second to get $Da\leq(1-d)r_1+br_2$. Multiply the second by
$1-\beta$ and add $c$ times the first to obtain the other bound.
\end{proof}

\begin{proposition}[Approximate evaluation and global validation]
Suppose only $\widetilde T(u)$ is evaluated, with verified block errors
$d_X(T_1(u),\widetilde T_1(u))\leq\delta_1$ and
$d_Y(T_2(u),\widetilde T_2(u))\leq\delta_2$. The preceding coordinate
bound remains valid after replacing $r_i$ by
$d_i(u_i,\widetilde T_i(u))+\delta_i$.
If a scalar residual function $r(s)$ is $L$-Lipschitz on a metric domain
$\mathcal D$ and $G$ is an $h$-net, then
\[
 \sup_{s\in\mathcal D}|r(s)|\leq\max_{g\in G}|r(g)|+Lh.
\]
The same statement holds for norm-valued residuals using their norm
Lipschitz bound. Truncation, quadrature and distribution-compression errors
must be included in the $\delta_i$ or in the grid bound.
\end{proposition}
\begin{proof}
The first claim is the triangle inequality for each block. For every $s$
choose $g\in G$ with distance at most $h$; then
$|r(s)|\leq |r(g)|+Lh$. Take the supremum. For vectors apply the reverse
triangle inequality to their norms.
\end{proof}

This validates all states only when the asserted net and Lipschitz bound
are themselves established. Finite simulation paths need not form a net
of a distribution space.

\section{Explicit failures of unjustified certificates}
\begin{proposition}[Conditioning cannot be omitted]
For each $\gamma\in(0,1)$ the scalar fixed-point equation
$x=(1-\gamma)x$ has unique equilibrium zero and residual
$r(x)=\gamma|x|$. Every error bound $|x|\leq C r(x)$ valid for nonzero
$x$ requires $C\geq\gamma^{-1}$. For the equation $F(x)=x^2=0$ on
$[0,1]$, no finite constant $C$ satisfies $|x|\leq C|F(x)|$ throughout
any neighborhood of zero. Finally fixed-price contraction of a household
block does not imply contraction of the coupled update.
\end{proposition}
\begin{proof}
The first ratio is identically $\gamma^{-1}$. The second ratio is $1/x$
for $x>0$, unbounded near zero. For the last assertion take
$T(v,p)=(\tfrac12v+p,v)$ on $\R^2$. For fixed $p$ the first component
contracts in $v$ with coefficient $1/2$. Its update matrix has positive
eigenvalue $(1+\sqrt{17})/4>1$, so iteration expands along its eigenvector
in every norm. Here the feedback product $bc=1$ violates $D>0$.
\end{proof}

Binding constraints do not invalidate the theorem when the update is
Lipschitz but nondifferentiable. They do invalidate applications that
silently replace a complementarity residual with an equality residual
on slack inequalities. No differentiability of $T$ was used.

\section{Propagation into estimation}
Let $\Theta$ be a parameter set with metric $d_\Theta$, true parameter
$\theta_0$, and model moments $M(\theta)\in\R^k$. Numerical moments
$\widetilde M$ and empirical moments $\widehat m$ satisfy
\[
 \sup_{\theta\in\Theta}\|\widetilde M(\theta)-M(\theta)\|\leq\eta,
 \quad \|\widehat m-M(\theta_0)\|\leq\zeta.
\]
Assume quantitative identification
$\|M(\theta)-M(\theta_0)\|\geq\kappa d_\Theta(\theta,\theta_0)$
for all $\theta$, with $\kappa>0$.

\begin{theorem}[Solver and sampling errors remain separately visible]
If $\widehat\theta$ is a $\xi$-approximate minimum-distance fit, meaning
\[
 \|\widetilde M(\widehat\theta)-\widehat m\|
 \leq\inf_\theta\|\widetilde M(\theta)-\widehat m\|+\xi,
\]
then
$d_\Theta(\widehat\theta,\theta_0)\leq
 [2\eta+2\zeta+\xi]/\kappa$.
If moments are $L_M$-Lipschitz in the equilibrium object and that object
has a uniform numerical error bound $\Delta$, with direct integration
error at most $\eta_0$, one may take $\eta=L_M\Delta+\eta_0$.
\end{theorem}
\begin{proof}
Use $\theta_0$ as a feasible comparison to bound the fitted discrepancy
by $\eta+\zeta+\xi$. A second triangle inequality bounds
$\|M(\widehat\theta)-M(\theta_0)\|$ by
$\eta+(\eta+\zeta+\xi)+\zeta$. Divide by $\kappa$.
The last assertion follows by adding the Lipschitz change in the moment
to its direct numerical integration error.
\end{proof}

\section{Verification and unresolved step}
All displayed bounds are deterministic and require no empirical
benchmark claims. In particular the exact inverse contains the feedback
denominator $D$, which can be small even when $\beta<1$ is well below one.
The first unresolved economic step is construction of an equilibrium-
equivalent invariant update with certified $b,c,d$ satisfying $D>0$ for a
nontrivial globally nonlinear aggregate-risk model across binding
constraints. The theorem supplies no general complexity guarantee for
covering its infinite-dimensional distribution domain. Thus the complete
catalogue computation and benchmarking programme remains unresolved.

\begin{thebibliography}{9}
\bibitem{ssj} A. Auclert, B. Bard\'oczy, M. Rognlie and L. Straub (2021),
Using the Sequence-Space Jacobian to Solve and Estimate Heterogeneous-Agent
Models, \emph{Econometrica} 89(5), 2375--2408.
\url{https://doi.org/10.3982/ECTA17434}.
\bibitem{ks} P. Krusell and A. A. Smith, Jr. (1998), Income and Wealth
Heterogeneity in the Macroeconomy, \emph{Journal of Political Economy}
106(5), 867--896. \url{https://doi.org/10.1086/250034}.
\end{thebibliography}

\section{Completion Estimate}
15\%

\end{document}
